\chapter{Tuberculosis}
\index{Tuberculosis}

\section*{Definition and Overview}
Tuberculosis (TB) is a chronic, infectious disease caused by bacteria belonging to the \textit{Mycobacterium tuberculosis} complex. Primarily affecting the lungs (pulmonary tuberculosis), it can also attack virtually any organ in the human body (extrapulmonary tuberculosis), including the lymph nodes, kidneys, brain, spine, and bones. Historically known by names such as ``consumption,'' ``phthisis,'' and the ``White Plague,'' tuberculosis has plagued humanity for millennia, with evidence of its existence found in Egyptian mummies dating back thousands of years. Despite being preventable and curable, TB remains one of the world's leading infectious killers, posing a persistent global public health challenge.

The significance of tuberculosis lies not only in its clinical severity but also in its complex socio-economic impact. It is intricately linked to poverty, malnutrition, overcrowded living conditions, and inadequate healthcare systems. The transmission of TB occurs through the air when an individual with active pulmonary disease coughs, sneezes, speaks, or sings, releasing tiny droplets containing the bacilli into the environment. When another person inhales these infectious aerosols, the bacteria can establish an infection. However, not everyone infected with \textit{Mycobacterium tuberculosis} develops active disease. In most individuals, the immune system successfully contains the bacteria, resulting in a state known as Latent Tuberculosis Infection (LTBI). Individuals with LTBI are asymptomatic, non-infectious, and may harbor the bacteria for decades without ever becoming ill. However, if the host's immune system becomes compromised due to aging, illness, or immunosuppressive therapy, the latent bacteria can reactivate, leading to active, symptomatic, and potentially infectious tuberculosis. The emergence of drug-resistant strains of the bacteria, particularly multidrug-resistant TB (MDR-TB) and extensively drug-resistant TB (XDR-TB), has further complicated global control efforts, threatening to dismantle decades of progress in combating this ancient disease.

\section*{Epidemiology}
Tuberculosis is a global pandemic, with its burden distributed unevenly across different geographic regions and demographic groups. According to the World Health Organization (WHO), TB is a leading cause of death from a single infectious agent, often ranking above HIV/AIDS. Globally, an estimated 10 million people contract active tuberculosis each year, and the disease claims approximately 1.3 to 1.5 million lives annually, including both HIV-negative and HIV-positive individuals.

The geographical distribution of tuberculosis is highly skewed, with the vast majority of cases occurring in low- and middle-income countries. The WHO Southeast Asia, Western Pacific, and African regions bear the highest burden, accounting for over 80\% of the global TB incidence. Countries such as India, Indonesia, China, the Philippines, Pakistan, Nigeria, Bangladesh, and the Democratic Republic of the Congo collectively represent more than two-thirds of the world's total TB cases.

Demographically, tuberculosis predominantly affects adults in their most productive years, typically between the ages of 15 and 59. However, no age group is immune, and pediatric tuberculosis represents a significant and often underdiagnosed segment of the epidemic, particularly in high-burden settings. Men are disproportionately affected by TB, accounting for approximately 56\% of all cases globally, compared to 33\% in adult women and 11\% in children. This gender disparity is thought to be driven by a combination of biological factors and socio-cultural behaviors, such as higher rates of smoking, alcohol consumption, and occupational exposure among men.

Certain populations are at an exceptionally high risk of contracting TB and experiencing poor treatment outcomes. People living with HIV (PLHIV) are 16 to 18 times more likely to develop active tuberculosis than those without HIV, as the virus destroys the cell-mediated immunity required to contain the mycobacteria. Consequently, TB is the leading cause of death among people living with HIV, accounting for about one-third of all AIDS-related deaths. Other highly vulnerable groups include prisoners, healthcare workers, refugees, migrants, individuals experiencing homelessness, and marginalized indigenous populations, who often face barriers to accessing timely diagnostic and treatment services.

\section*{Aetiology and Pathophysiology}
The primary causative agent of tuberculosis is \textit{Mycobacterium tuberculosis}, a slender, rod-shaped, non-motile, aerobic bacterium. It is characterized by a unique, lipid-rich cell wall containing high concentrations of mycolic acids. This waxy outer membrane gives the bacterium several distinctive properties, including acid-fastness (resistance to decolorization by acid during laboratory staining), resistance to many common antibiotics, and the ability to survive inside macrophages, the very cells designed to destroy invading pathogens.

The pathogenesis of tuberculosis is a complex interplay between the virulence factors of the mycobacterium and the host's cell-mediated immune response. The process unfolds through the following key stages and mechanisms:

\begin{itemize}
    \item \textbf{Inhalation and Alveolar Entry:} The infection begins when a susceptible host inhales microscopic droplet nuclei (1 to 5 micrometers in diameter) containing \textit{Mycobacterium tuberculosis} bacilli. These tiny droplets bypass the upper respiratory defenses and settle in the distal alveoli of the lungs.
    \item \textbf{Macrophage Phagocytosis and Survival:} Alveolar macrophages phagocytose the bacteria. Under normal circumstances, macrophages destroy ingested pathogens by fusing the phagosome (the vesicle containing the bacterium) with a lysosome filled with digestive enzymes and acidic compounds. However, \textit{M. tuberculosis} employs sophisticated survival mechanisms, including the inhibition of phagosome-lysosome fusion and the prevention of phagosomal acidification, allowing it to replicate intracellularly.
    \item \textbf{Primary Focus and Lymphatic Spread:} As the bacteria multiply within the macrophages, they eventually cause the host cell to rupture. The released bacilli infect neighboring cells and are carried by the lymphatic system to the regional lymph nodes (hilar lymph nodes). The combination of the initial pulmonary lesion (often in the middle or lower lobes) and the affected draining lymph node is known as the Ghon complex.
    \item \textbf{Granuloma Formation:} Approximately two to eight weeks after infection, a cell-mediated immune response is initiated. T-lymphocytes (specifically CD4+ helper T-cells) release cytokines, such as interferon-gamma (IFN-$\gamma$) and tumor necrosis factor-alpha (TNF-$\alpha$), which activate macrophages and recruit other immune cells to the site of infection. These cells organize to form a spherical structure called a granuloma (or tubercle). The granuloma consists of a central core of infected macrophages, multinucleated giant cells (Langhans cells), and necrotic debris, surrounded by a ring of lymphocytes and fibroblasts. This structure acts as a physical and immunological barrier, containing the bacteria and preventing further dissemination.
    \item \textbf{Caseous Necrosis:} Within the center of the mature granuloma, cells undergo a form of death known as caseous necrosis, characterized by a soft, cheese-like appearance. In most healthy individuals, the environment within the caseous center is acidic, hypoxic, and low in nutrients, which inhibits the growth of the mycobacteria, leading to the establishment of Latent Tuberculosis Infection (LTBI). Over time, these lesions may undergo calcification, visible on chest radiographs as Ranke complexes.
    \item \textbf{Reactivation and Liquefaction:} If the host's immune system becomes compromised (e.g., due to HIV infection, malnutrition, aging, or immunosuppressive drugs like TNF inhibitors), the cell-mediated control fails. The caseous material within the granuloma undergoes liquefaction, creating a rich medium that promotes rapid extracellular multiplication of the bacteria. The granuloma wall ruptures, spilling millions of bacilli into the bronchial tree. This leads to the formation of cavities in the lungs, facilitating the aerosol transmission of the bacteria to other individuals and causing systemic dissemination to other organs.
\end{itemize}

Several clinical risk factors significantly increase the likelihood of progressing from latent infection to active disease. These include:

\begin{itemize}
    \item \textbf{Immunosuppressive Conditions:} Human Immunodeficiency Virus (HIV) infection is the strongest risk factor for TB progression. Other conditions include diabetes mellitus, end-stage renal disease, silicosis, and certain cancers (such as leukemia and lymphoma).
    \item \textbf{Immunosuppressive Therapies:} Prolonged use of systemic corticosteroids, tumor necrosis factor-alpha (TNF-$\alpha$) antagonists, and drugs used to prevent organ transplant rejection severely impair the immune system's ability to maintain granulomas.
    \item \textbf{Substance Abuse and Malnutrition:} Chronic alcohol abuse, tobacco smoking, intravenous drug use, and severe malnutrition impair general immune function and damage pulmonary defense mechanisms.
    \item \textbf{Age:} Infants and children under the age of five, as well as the elderly, have immature or senescent immune systems, making them highly susceptible to rapid disease progression and severe forms of TB, such as miliary tuberculosis and tuberculous meningitis.
\end{itemize}

\section*{Clinical Features}
The clinical presentation of tuberculosis is highly variable and depends on whether the disease is pulmonary or extrapulmonary, the immune status of the host, and the age of the patient. Active pulmonary tuberculosis is the most common form, presenting with a mix of localized respiratory symptoms and systemic constitutional signs.

The classic clinical features of active pulmonary tuberculosis include:

\begin{itemize}
    \item \textbf{Chronic Cough:} The most common symptom is a persistent cough lasting for three weeks or longer. Initially, the cough may be dry and non-productive, but it typically becomes productive of mucoid, purulent, or blood-streaked sputum as the disease progresses.
    \item \textbf{Haemoptysis:} The coughing up of blood (haemoptysis) occurs when necrotizing lesions erode into bronchial blood vessels. It can range from minor blood streaking of sputum to life-threatening, massive hemorrhage.
    \item \textbf{Constitutional Symptoms:} These systemic signs of infection are mediated by the host's inflammatory response and include:
    \begin{itemize}
        \item \textbf{Unexplained Weight Loss and Anorexia:} A progressive loss of appetite and unintended weight reduction are highly characteristic.
        \item \textbf{Fever:} Patients often experience a low-grade, fluctuating fever, which frequently peaks in the afternoon or evening.
        \item \textbf{Drenching Night Sweats:} Profuse sweating during sleep is a classic symptom, often requiring patients to change their clothes or bedding.
        \item \textbf{Fatigue and Malaise:} A generalized, persistent sense of exhaustion, weakness, and ill health is common.
    \end{itemize}
    \item \textbf{Pleuritic Chest Pain:} Pain in the chest that worsens with breathing or coughing can occur if the tubercular infection involves the pleura (the membrane surrounding the lungs) or if a pleural effusion develops.
    \item \textbf{Dyspnoea:} Shortness of breath is usually a late feature, indicating extensive destruction of lung tissue, a large pleural effusion, or widespread miliary tuberculosis.
\end{itemize}

Extrapulmonary tuberculosis (EPTB) occurs when the bacteria disseminate from the lungs via the bloodstream or lymphatics to other organs. EPTB represents about 15\% to 20\% of cases in immunocompetent individuals but can account for more than 50\% of cases in patients co-infected with HIV. The clinical presentation of EPTB is organ-specific:

\begin{itemize}
    \item \textbf{Tuberculous Lymphadenitis (Scrofula):} This is the most common form of EPTB, typically affecting the cervical lymph nodes. It presents as painless, firm, slowly enlarging swellings in the neck that may eventually coalesce, fluctuate, and form chronic draining sinuses.
    \item \textbf{Pleural Tuberculosis:} Presents with fever, cough, and pleuritic chest pain, secondary to an inflammatory pleural effusion.
    \item \textbf{Skeletal Tuberculosis (Pott's Disease):} Involves the spine, typically the lower thoracic or lumbar vertebrae. It causes progressive back pain, spinal stiffness, and, if left untreated, can lead to vertebral collapse, kyphosis (humpback deformity), and neurological deficits due to spinal cord compression.
    \item \textbf{Tuberculous Meningitis:} The most severe form of EPTB, involving the central nervous system. It presents insidiously with headache, low-grade fever, neck stiffness, confusion, altered mental status, cranial nerve palsies, and can rapidly progress to coma and death if not treated urgently.
    \item \textbf{Genitourinary Tuberculosis:} Can affect the kidneys, ureters, bladder, and reproductive organs. Symptoms include sterile pyuria (pus in the urine with negative standard bacterial cultures), hematuria, dysuria, flank pain, and infertility.
    \item \textbf{Miliary Tuberculosis:} This is a life-threatening, disseminated form of TB resulting from the hematogenous spread of massive numbers of bacilli. It leads to the formation of tiny, millet-seed-like lesions throughout multiple organs, including the lungs, liver, spleen, and bone marrow. Clinically, it presents with severe constitutional symptoms, multi-organ dysfunction, and respiratory failure.
\end{itemize}

\section*{Differential Diagnosis}
Because the symptoms of tuberculosis---such as chronic cough, fever, weight loss, and night sweats---are non-specific, it is crucial to consider a wide range of alternative diagnoses. The differential diagnosis varies depending on the primary organ system involved:

\begin{itemize}
    \item \textbf{Community-Acquired and Atypical Pneumonia:} Bacterial infections (e.g., \textit{Streptococcus pneumoniae}, \textit{Mycoplasma pneumoniae}) and viral pneumonias can cause acute cough, fever, and pulmonary infiltrates. However, these typically have a more rapid onset and respond to standard antibiotic or supportive therapies.
    \item \textbf{Lung Cancer (Bronchogenic Carcinoma):} Particularly in older patients, smokers, or those with significant occupational exposure, lung cancer can mimic TB with chronic cough, haemoptysis, weight loss, and cavitary or nodular lesions on chest imaging.
    \item \textbf{Systemic Mycoses:} Fungal infections such as histoplasmosis, coccidioidomycosis, and blastomycosis can produce clinical and radiological findings virtually identical to TB, including chronic pulmonary symptoms, granuloma formation, and cavitation.
    \item \textbf{Non-Tuberculous Mycobacterial (NTM) Infections:} Infections caused by environmental mycobacteria (e.g., \textit{Mycobacterium avium} complex) can cause chronic pulmonary disease, especially in patients with pre-existing lung damage (like bronchiectasis or COPD) or immunocompromise.
    \item \textbf{Sarcoidosis:} A systemic inflammatory disease characterized by non-caseating granulomas. It often presents with bilateral hilar lymphadenopathy, pulmonary infiltrates, cough, and constitutional symptoms, making it highly challenging to distinguish from TB.
    \item \textbf{Lymphoma and Other Malignancies:} Can present with chronic fever, night sweats, weight loss, and generalized lymphadenopathy, closely mimicking tuberculous lymphadenitis or disseminated TB.
    \item \textbf{Inflammatory Bowel Disease (Crohn's Disease):} Gastrointestinal tuberculosis, which typically affects the ileocecal region, can present with abdominal pain, diarrhea, weight loss, and bowel wall thickening, highly resembling Crohn's disease.
\end{itemize}

\section*{When to Seek Medical Care}
Tuberculosis is a serious, progressive disease that requires prompt medical evaluation and treatment. Individuals should consult a medical professional if they experience any of the following symptoms:

\begin{itemize}
    \item A persistent cough that lasts for more than three weeks, especially if it produces thick, discolored, or blood-streaked sputum.
    \item Unexplained, progressive weight loss accompanied by a significant loss of appetite.
    \item Chronic, low-grade fevers, particularly when combined with profuse, drenching night sweats.
    \item Persistent swelling of lymph nodes, especially in the neck, axilla, or groin.
    \item Unexplained, worsening fatigue and generalized weakness that interferes with daily activities.
    \item Known exposure to an individual with active, infectious pulmonary tuberculosis, even in the absence of immediate symptoms.
\end{itemize}

Certain symptoms represent medical emergencies and require immediate transport to the nearest hospital emergency department or the contact of emergency services (such as **local emergency services** in many countries, **911** in the United States, or **999** in the United Kingdom):

\begin{itemize}
    \item \textbf{Severe Hemoptysis:} Coughing up large amounts of bright red blood (more than a few tablespoons), which indicates major vascular erosion and poses an immediate risk of airway obstruction and asphyxiation.
    \item \textbf{Acute Respiratory Distress:} Sudden or severe shortness of breath, rapid breathing, or a feeling of suffocation.
    \item \textbf{Neurological Changes:} Severe, sudden-onset headache accompanied by a stiff neck, high fever, confusion, altered consciousness, seizures, or focal neurological deficits (suggestive of tuberculous meningitis).
    \item \textbf{Severe, Crushing Chest Pain:} Which could indicate acute pericardial involvement (tuberculous pericarditis) or a cardiovascular event.
\end{itemize}

\section*{Diagnosis and Investigations}
The diagnosis of tuberculosis requires a multi-faceted approach combining clinical evaluation, immunological screening, radiological imaging, and microbiological confirmation. Accurate and timely diagnosis is critical to initiate therapy, prevent transmission, and curb the development of drug resistance.

Diagnostic investigations are broadly divided into screening tests for infection and confirmatory tests for active disease:

\begin{itemize}
    \item \textbf{Screening Tests for Infection (LTBI):} These tests detect the host's immune response to \textit{M. tuberculosis} antigens but cannot distinguish between active disease and latent infection.
    \begin{itemize}
        \item \textbf{Tuberculin Skin Test (TST) / Mantoux Test:} Involves the intradermal injection of a purified protein derivative (PPD) of \textit{M. tuberculosis} into the volar aspect of the forearm. The test is read 48 to 72 hours later by measuring the transverse diameter of induration (not erythema) in millimeters. The threshold for a positive result depends on the individual's risk factors (e.g., $\ge$5 mm for HIV-positive individuals, $\ge$10 mm for healthcare workers or immigrants from high-burden countries, $\ge$15 mm for individuals with no known risk factors). False positives can occur in individuals vaccinated with BCG (Bacillus Calmette-Gu\'{e}rin) or exposed to non-tuberculous mycobacteria.
        \item \textbf{Interferon-Gamma Release Assays (IGRAs):} In vitro blood tests (e.g., QuantiFERON-TB Gold, T-SPOT.TB) that measure the amount of interferon-gamma released by T-lymphocytes when mixed with specific \textit{M. tuberculosis} antigens. Unlike the TST, IGRAs do not cross-react with the BCG vaccine, making them highly specific, though they are more expensive and require specialized laboratory facilities.
    \end{itemize}
    \item \textbf{Radiological Imaging:}
    \begin{itemize}
        \item \textbf{Chest X-ray (CXR):} The primary imaging modality for suspected pulmonary TB. Classic findings of active post-primary TB include infiltrates in the apical and posterior segments of the upper lobes, segmentary or lobar consolidations, and cavitary lesions. In primary TB or immunocompromised individuals (like those with advanced HIV), findings may be atypical, presenting as middle or lower lobe infiltrates, hilar lymphadenopathy, or a normal chest radiograph.
        \item \textbf{Computed Tomography (CT) Scan:} Provides high-resolution detail of the lung parenchyma, mediastinal lymph nodes, and pleural space. It is particularly useful for identifying early cavitation, miliary patterns, bronchiectasis, and assessing extrapulmonary involvement.
    \end{itemize}
    \item \textbf{Microbiological Confirmation (The Gold Standard):}
    \begin{itemize}
        \item \textbf{Acid-Fast Bacilli (AFB) Smear Microscopy:} Sputum samples (ideally three consecutive morning specimens) are stained using the Ziehl-Neelsen or auramine-rhodamine method and examined under a microscope. While rapid, inexpensive, and useful for identifying the most infectious patients, it has low sensitivity (detecting only 50--60\% of cases) and cannot distinguish \textit{M. tuberculosis} from non-tuberculous mycobacteria.
        \item \textbf{Mycobacterial Culture:} Sputum or tissue samples are inoculated onto solid media (e.g., Lowenstein-Jensen) or liquid media (e.g., Mycobacterial Growth Indicator Tube - MGIT). Culturing is highly sensitive and allows for subsequent Drug Susceptibility Testing (DST). However, because \textit{M. tuberculosis} is a slow-growing organism, solid culture can take 3 to 8 weeks to yield results (liquid culture is faster, taking 10 to 21 days).
        \item \textbf{Nucleic Acid Amplification Tests (NAATs):} Rapid molecular tests, such as the GeneXpert MTB/RIF assay, detect \textit{M. tuberculosis} DNA directly from sputum samples within two hours. Importantly, these tests also detect mutations associated with resistance to rifampicin, a key first-line drug. The WHO recommends NAATs as the initial diagnostic test for all individuals with suspected pulmonary TB.
    \end{itemize}
    \item \textbf{Pathology (Biopsy):} Histopathological examination of affected tissues (e.g., lymph nodes, pleura, liver) reveals characteristic caseating granulomas, Langhans multinucleated giant cells, and acid-fast bacilli, supporting a presumptive diagnosis when microbiological specimens are difficult to obtain.
\end{itemize}

\section*{Management}
The primary goals of tuberculosis management are to cure the individual patient, minimize the risk of death and long-term complications, prevent relapse, and halt the transmission of the infection to others. The cornerstone of TB treatment is a prolonged course of combination antimicrobial therapy. Monotherapy must never be used to treat active TB, as it rapidly leads to the selection of drug-resistant mutants.

\subsection*{Treatment of Drug-Susceptible Pulmonary Tuberculosis}
For patients with confirmed or highly suspected drug-susceptible pulmonary TB, the standard regimen recommended by the WHO and international guidelines consists of a six-month, two-phase regimen:

\begin{itemize}
    \item \textbf{Intensive Phase (2 Months):} Utilizes four first-line bactericidal drugs to rapidly reduce the bacterial load, clinical symptoms, and infectivity.
    \begin{itemize}
        \item \textbf{Isoniazid (H):} Highly bactericidal, targeting cell wall synthesis (mycolic acid).
        \item \textbf{Rifampicin (R):} Potent bactericidal and sterilizing agent, inhibiting bacterial RNA synthesis.
        \item \textbf{Pyrazinamide (Z):} Active in acidic environments within macrophages, killing semi-dormant bacilli.
        \item \textbf{Ethambutol (E):} Bacteriostatic agent, inhibiting cell wall synthesis; added to prevent the emergence of drug resistance.
    \end{itemize}
    \item \textbf{Continuation Phase (4 Months):} Utilizes two drugs to eliminate remaining dormant or persisting bacilli, preventing relapse.
    \begin{itemize}
        \item \textbf{Isoniazid (H) and Rifampicin (R)} administered daily.
    \end{itemize}
\end{itemize}

To simplify administration and improve adherence, these medications are frequently combined into Fixed-Dose Combination (FDC) tablets.

\subsection*{Directly Observed Therapy (DOT)}
Adherence to the full, uninterrupted course of TB therapy is the single most critical factor in achieving a cure and preventing the development of drug-resistant strains. To support patients, public health programs utilize Directly Observed Therapy (DOT), where a trained healthcare provider or community worker visualizes the patient swallowing each dose of medication. DOT may be delivered in clinics, at home, or virtually via video-enabled technology (vDOT).

\subsection*{Management of Drug-Resistant Tuberculosis}
Drug-resistant TB is a major global threat. Multidrug-resistant TB (MDR-TB) is defined as resistance to at least isoniazid and rifampicin. Extensively drug-resistant TB (XDR-TB) involves resistance to isoniazid, rifampicin, any fluoroquinolone, and at least one second-line injectable drug.

Treatment of drug-resistant TB is highly complex, expensive, and associated with more severe side effects:
\begin{itemize}
    \item It requires longer treatment regimens, typically lasting 9 to 20 months depending on the resistance profile and the regimen chosen.
    \item It utilizes second-line agents, including fluoroquinolones (e.g., levofloxacin, moxifloxacin), bedaquiline, linezolid, pretomanid, and cycloserine.
    \item Treatment must be closely guided by specialists, using individualized regimens based on rapid molecular drug-susceptibility testing.
\end{itemize}

\subsection*{Monitoring and Adverse Effects}
Patients undergoing TB treatment must be monitored regularly for clinical response, adherence, and medication toxicity. Routine blood tests, including liver function tests (LFTs) and complete blood counts, are performed, as many first-line drugs are hepatotoxic. Key adverse effects include:
\begin{itemize}
    \item \textbf{Hepatotoxicity (Drug-Induced Liver Injury):} Caused by isoniazid, rifampicin, and pyrazinamide. Characterized by nausea, vomiting, abdominal pain, jaundice, and elevated liver enzymes.
    \item \textbf{Peripheral Neuropathy:} Caused by isoniazid, due to interference with pyridoxine (vitamin B6) metabolism. Prevented by co-administering daily pyridoxine supplements, especially in high-risk patients (pregnant women, elderly, malnourished, HIV-positive).
    \item \textbf{Optic Neuritis:} Caused by ethambutol. Presents as decreased visual acuity or impaired red-green color vision. Visual testing should be performed before and during therapy.
    \item \textbf{Hyperuricemia and Gout:} Caused by pyrazinamide, which inhibits renal excretion of uric acid.
    \item \textbf{Red-Orange Discoloration of Bodily Fluids:} A benign side effect of rifampicin, which colors urine, sweat, tears, and saliva.
\end{itemize}

\subsection*{Adjuvant Corticosteroids}
Systemic corticosteroids (e.g., prednisolone or dexamethasone) are indicated as adjunctive therapy in specific extrapulmonary manifestations, such as tuberculous meningitis and tuberculous pericarditis, to reduce intense local inflammation and improve survival rates.

\section*{Complications}
Untreated, delayed, or inadequately treated tuberculosis can lead to severe, permanent damage to the lungs and other organ systems, resulting in chronic disability or death.

Potential complications include:

\begin{itemize}
    \item \textbf{Permanent Lung Damage:} Chronic inflammation and cavitation can lead to extensive pulmonary fibrosis, bronchiectasis (permanent widening of the airways), and chronic obstructive pulmonary disease (COPD), resulting in persistent shortness of breath, recurrent respiratory infections, and respiratory failure.
    \item \textbf{Hemoptysis and Rasmussen's Aneurysm:} Severe or recurrent bleeding can occur. A Rasmussen's aneurysm is a dilation of a pulmonary artery branch within a tuberculous cavity, which can rupture and cause fatal, massive hemoptysis.
    \item \textbf{Secondary Bacterial Infections:} Cavities left behind by cured TB can become colonized by environmental fungi, particularly \textit{Aspergillus} species, forming a fungal ball (aspergilloma) that causes recurrent bleeding.
    \item \textbf{Pneumothorax:} Rupture of a subpleural cavity or caseous lesion into the pleural space can lead to the collapse of the lung (pneumothorax) or a collection of pus (empyema).
    \item \textbf{Spinal Deformity and Paralysis:} Untreated skeletal TB (Pott's disease) can cause collapse of the vertebral bodies, leading to severe kyphosis, spinal instability, and compression of the spinal cord resulting in paraplegia.
    \item \textbf{Neurological Deficits:} Tuberculous meningitis can lead to hydrocephalus, cranial nerve palsies, cognitive impairment, stroke, epilepsy, and permanent motor deficits.
    \item \textbf{Adrenal Insufficiency (Addison's Disease):} Tuberculous destruction of the adrenal glands is a classic cause of primary adrenal insufficiency, presenting with chronic fatigue, low blood pressure, electrolyte abnormalities, and life-threatening adrenal crisis.
    \item \textbf{Infertility:} Genitourinary tuberculosis can cause scarring and obstruction of the fallopian tubes in women (leading to pelvic inflammatory disease and infertility) and the epididymis or seminal vesicles in men.
    \item \textbf{Constrictive Pericarditis:} Tuberculous involvement of the pericardium can lead to dense fibrosis, calcification, and constriction of the heart, impairing cardiac filling and causing chronic right-sided heart failure.
\end{itemize}

\section*{Prognosis and Outcomes}
The prognosis of tuberculosis is highly favorable if the disease is diagnosed early and treated with an appropriate, fully completed regimen of antimicrobial therapy. In patients with drug-susceptible pulmonary tuberculosis who adhere to their treatment, cure rates exceed 95\%.

However, the prognosis worsens significantly under several circumstances:
\begin{itemize}
    \item \textbf{Co-infection with HIV:} HIV-positive patients have a much higher mortality rate, particularly if they do not receive antiretroviral therapy (ART) concurrently with TB treatment. ART significantly reduces mortality in this population.
    \item \textbf{Drug Resistance:} Patients with MDR-TB and XDR-TB face substantially lower cure rates, typically ranging from 50\% to 70\% for MDR-TB, and even lower for XDR-TB. The treatment is longer, less effective, and associated with higher rates of treatment failure, relapse, and death.
    \item \textbf{Delayed Diagnosis:} Delays in initiating treatment allow the disease to progress, increasing the risk of severe complications, irreversible tissue damage, and transmission.
    \item \textbf{Extrapulmonary Involvement:} Certain forms of EPTB, particularly tuberculous meningitis and miliary tuberculosis, carry a high risk of mortality (up to 20--50\%, even with treatment) and permanent neurological or functional sequelae.
    \item \textbf{Patient Comorbidities:} Advanced age, severe malnutrition, chronic liver or renal disease, and active substance abuse are associated with higher rates of drug toxicity, treatment discontinuation, and mortality.
\end{itemize}

Without treatment, active tuberculosis is progressive and often fatal. Historical studies indicate that approximately 50\% to 70\% of individuals with untreated, sputum-smear-positive pulmonary TB die within 5 to 10 years, with the majority dying within 2 years of onset.

\section*{Prevention and Self-Care}
Preventing the spread of tuberculosis requires a combination of clinical interventions, public health policies, environmental controls, and personal self-care measures. Because TB is airborne, community-level transmission control is paramount.

Key prevention and self-care strategies include:

\begin{itemize}
    \item \textbf{Vaccination (BCG Vaccine):} The Bacillus Calmette-Gu\'{e}rin (BCG) vaccine is a live-attenuated strain of \textit{Mycobacterium bovis}. It is widely administered to infants at birth in countries with a high prevalence of TB. While its efficacy in preventing adult pulmonary TB is variable and limited, BCG is highly effective (up to 80\%) in protecting infants and young children against severe, disseminated forms of the disease, including tuberculous meningitis and miliary TB. It is generally not recommended for routine use in low-burden countries, except for specific high-risk individuals.
    \item \textbf{Treatment of Latent TB Infection (LTBI):} Identifying and treating individuals with latent infection (preventative therapy) is a crucial pillar of TB elimination strategies, particularly in low-incidence settings. Standard regimens include daily isoniazid for 6 to 9 months, or shorter, newer regimens such as weekly isoniazid and rifapentine for 3 months, or daily rifampicin for 4 months.
    \item \textbf{Infection Control Measures:}
    \begin{itemize}
        \item \textbf{Administrative Controls:} Rapid triaging, isolation, and initiation of treatment for suspected or confirmed active TB cases in healthcare facilities.
        \item \textbf{Environmental Controls:} Utilizing negative pressure isolation rooms, high-efficiency particulate air (HEPA) filtration systems, and ultraviolet germicidal irradiation (UVGI) to clear the air of infectious droplets.
        \item \textbf{Personal Protective Equipment (PPE):} Healthcare workers caring for infectious patients must wear fit-tested N95 or higher-level respirators.
    \end{itemize}
    \item \textbf{Self-Care for Infectious Patients:}
    \begin{itemize}
        \item \textbf{Isolation:} Patients with active, infectious pulmonary TB must remain isolated at home during the initial weeks of treatment (usually until they have completed at least two weeks of effective therapy and have consecutive negative sputum smears).
        \item \textbf{Cough Etiquette:} Always cover the mouth and nose with a tissue when coughing, sneezing, or laughing, and dispose of the tissue immediately in a closed bin. Alternatively, cough or sneeze into the inner elbow.
        \item \textbf{Ventilation:} Keep windows open as much as possible to allow fresh air to circulate and dilute any infectious aerosols. Avoid sleeping in the same room as uninfected family members.
        \item \textbf{Avoid Public Spaces:} Do not go to work, school, or use public transport until cleared by a physician or public health authority.
    \end{itemize}
    \item \textbf{Adherence Support:} Patients must take their medications exactly as prescribed, without skipping doses or stopping early, even if they begin to feel better. Setting reminders, using pill organizers, and maintaining open communication with the healthcare team are essential.
    \item \textbf{Nutritional and Lifestyle Support:} A balanced diet rich in protein and micronutrients is vital to support the immune system and repair tissue damage. Smoking cessation and avoiding alcohol are critical, as they improve lung function and reduce the risk of drug-induced liver injury.
\end{itemize}

\section*{Sources and Further Reading}
\begin{itemize}
    \item \textbf{World Health Organization (WHO) - Tuberculosis:} Comprehensive global data, guidelines, policy briefs, and annual global TB reports. URL: \texttt{https://www.who.int/health-topics/tuberculosis}
    \item \textbf{Centers for Disease Control and Prevention (CDC) - Tuberculosis (TB):} Extensive clinical resources, patient education materials, and guidelines for testing and treatment. URL: \texttt{https://www.cdc.gov/tb/index.html}
    \item \textbf{The Union (International Union Against Tuberculosis and Lung Disease):} A leading global scientific organization dedicated to ending tuberculosis through research and education. URL: \texttt{https://theunion.org}
    \item \textbf{Stop TB Partnership:} A global initiative hosting a wide range of partners working to accelerate progress towards TB elimination. URL: \texttt{https://www.stoptb.org}
    \item \textbf{Mayo Clinic - Tuberculosis:} A patient-oriented medical reference outlining symptoms, causes, diagnosis, and treatment options. URL: \texttt{https://www.mayoclinic.org/diseases-conditions/tuberculosis/symptoms-causes/syc-20351250}
    \item \textbf{National Health Service (NHS) - Tuberculosis (TB):} A clear, accessible guide to TB symptoms, treatment, and prevention for patients in the United Kingdom. URL: \texttt{https://www.nhs.uk/conditions/tuberculosis-tb/}
\end{itemize}